% Part: methods
% Chapter: induction
% Section: structural-induction

\documentclass[../../../include/open-logic-section]{subfiles}

\begin{document}

\olfileid{mth}{ind}{sti}

\olsection{संरचनात्मक विगमन}

आत्तापर्यंत सर्व नैसर्गिक संख्यांविषयी निष्कर्ष प्रस्थापित करण्यासाठी विगमन वापरले. पण त्याच्याशी संबंधित तत्त्व थेट वापरून विगमनाधारित व्याख्येने दिलेल्या संचातील सर्व !!{element}s विषयी निष्कर्ष सिद्ध करता येतात. याला अनेकदा \emph{संरचनात्मक} विगमन म्हणतात; कारण ते त्या व्याख्येने दिलेल्या वस्तूंच्या संरचनेवर अवलंबून असते.

सहसा विगमनाधारित व्याख्येत (a) संचातील “आरंभीच्या” !!{element}s ची यादी आणि (b) जुन्यांपासून नवे !!{element}s निर्माण करणाऱ्या क्रियांची यादी दिलेली असते. उदाहरणार्थ, नीटस पदांच्या प्रकरणात आरंभीच्या वस्तू म्हणजे अक्षरे. एकच क्रिया आहे:
\begin{align*}
  o(s_1, s_2) = & [s_1 \circ s_2]
\end{align*}
नैसर्गिक संख्यांचा संच~$\Nat$ देखील विगमनाधारित व्याख्येने दिलेला आहे असा विचार करता येतो: आरंभीची वस्तू~$0$ आणि क्रिया म्हणजे~$x + 1$ हे उत्तरवर्ती फलन.

विगमनाधारित व्याख्येने दिलेल्या संचातील सर्व घटकांविषयी काही सिद्ध करायचे असेल---म्हणजे संचातील प्रत्येक !!{element} मध्ये~$P$ हा गुणधर्म आहे हे सिद्ध करायचे असेल---तर पुढील गोष्टी कराव्या लागतात:
\begin{enumerate}
\item आरंभीच्या वस्तूंमध्ये~$P$ आहे हे सिद्ध करा.
\item प्रत्येक~$o$ क्रियेसाठी, तिच्या सर्व वितर्कांत~$P$ असेल, तर फलितातही तो असतो हे सिद्ध करा.
\end{enumerate}
उदाहरणार्थ, सर्व नीटस पदांविषयी काही सिद्ध करण्यासाठी ते सर्व अक्षरांविषयी सत्य आहे हे सिद्ध करू; तसेच $s_1$ आणि $s_2$ यांविषयी ते स्वतंत्रपणे सत्य असेल, तर $[s_1 \circ s_2]$ विषयीही सत्य असते हे सिद्ध करू.

\begin{prop}
कोणत्याही नीटस पदात~$t$ मध्ये $[$ ची संख्या $]$ च्या संख्येइतकी असते.
\end{prop}

\begin{proof}
संरचनात्मक विगमन वापरू. नीटस पदांची व्याख्या विगमनाधारित आहे: अक्षरे या आरंभीच्या वस्तू आणि जुन्या पदांपासून नवी नीटस पदे रचण्यासाठी $o$ ही क्रिया.
\begin{enumerate}
\item प्रत्येक अक्षरासाठी विधान सत्य आहे; कारण स्वतंत्र अक्षरात $[$ ची संख्या~$0$ आणि $]$ ची संख्याही~$0$ आहे.
\item $s_1$ मध्ये $[$ ची संख्या $]$ च्या संख्येइतकी आहे आणि $s_2$ साठीही तेच सत्य आहे असे समजा. $o(s_1, s_2)$ मध्ये---म्हणजे $[s_1 \circ s_2]$ मध्ये---$[$ ची संख्या ही $s_1$ आणि $s_2$ यांमधील $[$ च्या संख्यांची बेरीज अधिक एक आहे. $o(s_1, s_2)$ मध्ये $]$ ची संख्या ही $s_1$ आणि $s_2$ यांमधील $]$ च्या संख्यांची बेरीज अधिक एक आहे. म्हणून $o(s_1, s_2)$ मधील $[$ ची संख्या $o(s_1,s_2)$ मधील $]$ च्या संख्येइतकी आहे.
\end{enumerate}
\end{proof}

\begin{prob}
कोणत्याही नीटस पदाची सुरुवात~$]$ ने होत नाही, हे संरचनात्मक विगमनाने सिद्ध करा.
\end{prob}

संरचनात्मक विगमनाने आणखी एक सिद्धता देऊ. चिन्हांच्या~$t$ चिन्हमालेचा अरिक्त उचित पूर्वखंड म्हणजे अशी अरिक्त~$s$ चिन्हमाला की डावीकडून वाचताना प्रत्येक चिन्हानुसार ती $t$ शी जुळते, पण $t$ अधिक लांब असते. उदाहरणार्थ, $[a \circ {}$ हा $[a \circ b]$ चा अरिक्त उचित पूर्वखंड आहे. मात्र $[b \circ {}$ तसा नाही---दुसऱ्या चिन्हावर फरक आहे---आणि $[a \circ b]$ देखील तसा नाही---दोन्हींची लांबी समान आहे. येथे अरिक्त पूर्वखंडच घेतले आहेत: रिकाम्या पूर्वखंडात उघडणारे आणि बंद करणारे दोन्ही कंस शून्य असल्यामुळे पुढील काटेकोर असमानता त्यासाठी खरी नाही.

\begin{prop}\ollabel{prop:initial}
नीटस पद~$t$ च्या प्रत्येक अरिक्त उचित पूर्वखंडात $[$ ची संख्या $]$ पेक्षा जास्त असते.
\end{prop}

\begin{proof}
~$t$ वर विगमन करू:
\begin{enumerate}
\item $t$ हे स्वतंत्र अक्षर आहे: मग $t$ ला कोणतेही अरिक्त उचित पूर्वखंड नाहीत.
\item $t = [s_1 \circ s_2]$, जिथे $s_1$ आणि~$s_2$ ही नीटस पदे आहेत. $r$ हा $t$ चा अरिक्त उचित पूर्वखंड असेल, तर काही शक्यता आहेत:
\begin{enumerate}
\item $r$ हे केवळ $[$ आहे: मग $r$ मध्ये $[$ ची संख्या~$]$ पेक्षा एकाने जास्त आहे.
\item $r$ हे $[r_1$ आहे, जिथे $r_1$ हा~$s_1$ चा अरिक्त उचित पूर्वखंड आहे. $s_1$ नीटस पद असल्यामुळे विगमन गृहीतकानुसार $r_1$ मध्ये $[$ हे $]$ पेक्षा जास्त आहेत; आणि $[r_1$ साठीही तेच खरे आहे.
\item $r$ हे $[s_1$ किंवा $[s_1 \circ {}$ आहे. आधीच्या निष्कर्षानुसार~$s_1$ मध्ये $[$ आणि $]$ यांच्या संख्या समान आहेत. म्हणून $[s_1$ किंवा $[s_1 \circ {}$ मध्ये $[$ ची संख्या~$]$ पेक्षा एकाने जास्त आहे.
\item $r$ हे $[s_1 \circ r_2$ आहे, जिथे $r_2$ हा~$s_2$ चा अरिक्त उचित पूर्वखंड आहे. विगमन गृहीतकानुसार $r_2$ मध्ये $[$ हे $]$ पेक्षा जास्त आहेत. आधीच्या निष्कर्षानुसार~$s_1$ मध्ये $[$ आणि $]$ यांच्या संख्या समान आहेत. म्हणून $[s_1 \circ r_2$ मध्ये $[$ ची संख्या~$]$ पेक्षा मोठी आहे.
\item $r$ हे $[s_1 \circ s_2$ आहे. आधीच्या निष्कर्षानुसार $s_1$ मध्ये $[$ आणि $]$ यांच्या संख्या समान आहेत आणि~$s_2$ मध्येही त्या समान आहेत. म्हणून $[s_1 \circ s_2$ मध्ये $[$ ची संख्या~$]$ पेक्षा एकाने जास्त आहे.
\end{enumerate}
\end{enumerate}
\end{proof}

\end{document}
